\chapter{Conclusiones}
\thispagestyle{fancy}
\label{ch:conclusiones}

Esta práctica logró integrar los conceptos y modelos teóricos de la cinemática,
donde los tres tipos de movimientos rectilíneos (uniformemente acelerado, 
uniforme y uniformemente retardado) se analizaron en un sistema de tres tramos, 
representando el recorrido de un autobús y la propagación de incertidumbres en cada fase.
Se determinó que el autobús recorre una distancia total de $(1320{,}0 \pm 8{,}85)$~m en un tiempo total de
$(60{,}00 \pm 0{,}03)$~s, con incertidumbres relativas inferiores al $1\%$ en
ambos casos.

Se identificó que la incertidumbre dominante proviene del Tramo~II, donde el
error en la velocidad $\Delta V_1 = 0{,}16$~m/s se amplifica por el largo
intervalo de tiempo de 50~s, generando $\Delta X_2 = 8{,}24$~m. A pesar de que 
el error relativo es bajo, su impacto absoluto se encuentra significativo
en comparación con las incertidumbres en otros tramos. Para reducir 
el impacto de esta propagación de error y mejorar la precisión de los resultados, se
recomienda mejorar la precisión en la determinación de 
la aceleración del Tramo~I, ya que esta propaga su error hacia $V_1$ y 
de allí hacia toda la trayectoria de velocidad constante. 

Las representaciones gráficas de posición vs.\ tiempo confirmaron visualmente
la naturaleza de cada tipo de movimiento y los ajustes polinomiales devolvieron
$R^2 \approx 1{,}000$ en los tres tramos; un resultado que confirma la certeza 
de los modelos teóricos aplicados, La herramienta computacional
utilizada demostró ser equivalente en sus resultados, lo que añade confianza a
los cálculos realizados.

En conjunto con el Ejercicio~A se logró identificar dos factores dominantes
 en la propagación de incertidumbre
al momento de realizar las mediciones: la reproducibilidad de las condiciones 
experimentales y la precisión de los instrumentos de medición. Sin embargo, 
también se destaca la importancia de realizar múltiples mediciones y 
promediar los resultados para reducir el impacto de errores aleatorios.